\noindent\textbf{Checked core.} The declaration \texttt{ADMMCounterexample.counterexample} compiled with exit code zero. It contains $1<\tau<2$, convexity, compactness, and closedness of both sets, the common feasible origin, a zero fixed trajectory, every update of the displayed trajectory, five-periodicity, nonconstancy, the uniform squared-residual bound, and the assertion that this residual does not tend to zero. The actual projection minimizers are checked over the full real convex hulls, rather than only at their vertices.

The three core modules have no \texttt{sorry}, \texttt{admit}, custom \texttt{axiom}, \texttt{native\_decide}, or externally trusted numerical certificate. The final axiom report is exactly \texttt{[propext, Classical.choice, Quot.sound]}. These are the standard logical foundations, not assumptions of the counterexample. The use of \texttt{noncomputable} for real arithmetic changes executable-code generation, not the proof obligation or the kernel's checking of the theorem.

The indicator-objective and saddle-point interpretation is proved in the text. The affine determinants, ordinary interiority, the parameter interval, and analytic $C^2$/strong-convexity/Lipschitz-gradient properties of the extensions are not part of this core Lean declaration. Their complete mathematical arguments and exact arithmetic certificates are supplied separately; no formal-verification claim is made for those additional regularity statements.

\noindent\textbf{Checked finite-objective extension.} The separate declaration \path{ADMMHingeTrajectory.verified_orbit_and_optimizer} also compiled successfully, with the same three standard foundational dependencies. For the actual functions in \eqref{eq:smoothf}--\eqref{eq:smoothg}, it certifies global proximal minimizers at every iteration, the multiplier recurrence, the unique constrained optimizer zero, periodicity, nonconstancy, and the failure of the residual to tend to zero. The analytic regularity statements about these functions are established by the handwritten derivative and Hessian calculations above; those properties are not encoded in this additional Lean theorem.
